% THETA3-A6 contribution by /root/theta_032, 21 November 2026.
% GPT-6 Astra Pro; inherited reasoning-effort label not separately exposed.
% Liu/BMM/CM classification are cited inputs, not independently re-proved.

\subsection{THETA3-A6: a direct construction for balanced tensors of arbitrary even length}\label{A6}

\paragraph{Statement and scope.}
Let $n\geq2$, let $E$ be a Galois CM field with real subfield $F$
and $[F:\Q]\geq2$, and let $c\in\operatorname{Gal}(E/\Q)$ be complex
conjugation. Suppose that the odd sign functions
$v_1,\ldots,v_{2n}:\operatorname{Gal}(E/\Q)\to\{1,-1\}$ satisfy
\begin{equation}\label{A6-balance}
 v_i(1)=1,\qquad \sum_{i=1}^n v_i=\sum_{i=n+1}^{2n}v_i.
\end{equation}
Put $\lambda=(1,\ldots,1,c,\ldots,c)$, with $n$ entries of each kind.
The theta construction on a compact ball quotient of dimension $n$
gives a direct cycle construction for the line
\begin{equation}\label{A6-line}
 \bigotimes_{i=1}^{2n}U(v_i)_{\lambda_i}
 \ \subset\ H^{2n}\!\left(\prod_{i=1}^{2n}B(v_i),\overline\Q\right).
\end{equation}
The arithmetic input is Liu's exact oscillator-component theorem in
arbitrary rank, together with the reflex-type argument of
van Hoften--Sempliner. The proof below does not infer the prescribed
$U(v_i)$ sources from CM-ness of the Albanese alone. It does not use
the reduction to four-factor relations.

All auxiliary abelian varieties in this construction remain CM. The algebraicity
conclusion is already covered by the supplied CM Hodge-conjecture
results; the contribution here is the direct $2n$-factor construction,
its exact multiplicity and differential degree, and the comparison
with the stated theta inputs. No claim of research priority is made.

\paragraph{The exact arithmetic source in rank $m\geq3$.}
Let $V$ be Hermitian of rank $m$, with signature $(m-1,1)$ at the
distinguished real place and $(m,0)$ elsewhere. It is anisotropic
over $E$, because there is a definite real place. Replacing its active
real component by a positive definite component while retaining the
finite components gives a totally definite incoherent Hermitian
space. The original $V$ is its active nearby coherent space.
Consequently Liu's setup applies to this precise Shimura tower:
Definitions C.3--C.4 and Proposition C.5, printed pp.108--109, and
Section 4.2, pp.45--47. The compact case holds because $[F:\Q]>1$.
There is no all-place indefinite-signature hypothesis in this use
of Liu's theorem.

For an odd $v$ with $v(1)=1$, write
$\Phi(v)=\{g:v(g)=1\}$ and $d(\tau)=v(\tau^{-1})$. Choose a
unitary conjugate-symplectic Hecke character $\mu$ of weight one with
\[
 \mu|_{\A_F^\times}=\epsilon_{E/F},\qquad
 \Phi_\mu=\Phi(v)^{-1}.
\]
Existence is van Hoften--Sempliner, Lemma 3.1.5, pp.9--10.
Liu's admissibility condition is the existence of a skew element
$e\in E^\times$, $e^c=-e$, such that
$\operatorname{Im}\tau(e)<0$ for $\tau\in\Phi_\mu$, with the
specified finite norm classes; see Definition 4.12, p.47. This is
exactly the 032 sign prescription: for $e=\delta h_P$, it says
\[
 \operatorname{sgn}(h_{P,\tau})
   =-d(\tau)\operatorname{sgn}(\operatorname{Im}\tau(\delta)).
\]
The skew-Hermitian source line therefore has active signature
$(0,1)$.

Liu, Definition 4.3, p.41, puts
$\mu_{\rm alg}=\mu|\cdot|_E^{-1/2}$ and defines a coefficient
field $M_\mu$ containing the reflex field $M'_\mu$ of
$(E,\Phi_\mu)$, with reflex type $\Psi_\mu$.
Definition 4.5, p.42, specifies a CM abelian variety $A_\mu$ and its
distinguished $M_\mu$-inclusion eigenline, whose Galois character is
$\mu_{\rm alg}$; Proposition 4.6, pp.42--43, supplies such data.
Theorem 4.15, pp.50--51, identifies the Galois character of the
particular admissible oscillator component with this exact character
for every rank $m\geq3$. Theorem 4.18, pp.52--54, then gives
\begin{equation}\label{A6-Liu}
 \Hom_E(A_\infty,A_\mu)_\Q\otimes_{M_\mu}\C
       \simeq\bigoplus_{\varepsilon,\chi}
                         \omega(\mu,\varepsilon,\chi).
\end{equation}
The sum uses $\mu$-admissible $\varepsilon$. The map in its proof
is pullback of a basis of the distinguished inclusion eigenline;
its finite-level assertion identifies $K$-invariants with
$\Hom_E(A_K,A_\mu)_\Q$. Here $\chi$ is the finite central
automorphic character of Liu, Definition 4.11. It is neither the
splitting character $\mu$ nor the source torus character.
Corollary 4.20, pp.54--55, describes the whole CM Albanese, but the
stronger componentwise statement \eqref{A6-Liu} is what is needed.

Let $(E_0,\Phi_0)$ be the primitive CM type inducing
$(E,\Phi(v))$. Then $B(v)\sim B_0^{[E:E_0]}$. Since $E$ is
Galois and $\Phi_\mu=\Phi(v)^{-1}$, the reflex pair
$(M'_\mu,\Psi_\mu)$ is $(E_0,\Phi_0)$, and
$A_{\mu,\C}\sim B_0^{[M_\mu:E_0]}$. These are precisely the
CM-type computations in van Hoften--Sempliner, Corollary 3.1.13,
p.12, and Proposition 4.1.8, pp.15--16. Their proof with the
primitive factor and the selected inclusion line gives
\begin{equation}\label{A6-extraction}
 \omega(\mu,\varepsilon,\chi)
 \subseteq
 \operatorname{span}_{\C}\{h(U(v)_1):
       h\in\Hom_{\Hdg}(U(v),H^1(X,\Q))\},
\end{equation}
on a sufficiently deep common cover. This argument is independent
of $m$. The inverse in $\Phi(v)^{-1}$ is essential.

\paragraph{The first theta vector and the scalar conventions.}
The local comparison needed to use \eqref{A6-extraction} is explicit.
In a graph frame at a base point, 032's first Taylor tensor is the
ordinary gradient: a nonzero scalar times a linear position function
multiplied by the positive Gaussian. The Bargmann transform sends
linear position times the vacuum to a nonzero scalar times the
corresponding linear Fock coordinate. After projection to the active
cotangent summand, this is the contraction
\[
 \psi_{1,0}=\sum_{\alpha=1}^{m-1}
                 z''_\alpha\otimes\xi'_{\alpha,m}.
\]
It is the $q=1$ case of BMM, published Lemma 5.5, pp.43--44;
equation (5.11), p.54, defines its Schwartz realization by the
inverse Bargmann transform. One may check the comparison on this
fibre and transport it by $\SU(V)$, since both sides have the same
covariance. The Gaussian creation-operator formulas in BMM,
Appendix C, Lemmas C.2--C.3, pp.115--116, also verify directly that
the coefficient is nonzero.

The numerical scalar depends on the Bargmann, Gaussian and
constant-line normalizations. It need not be algebraic: only a
complex span and nonvanishing are used in this argument. It must be
retained if a numerical period identity is later asserted.
BMM, Lemma 5.18, p.49, fixes the twist
$\det^{(a-b)/2}$ required for the special cochain on $\U(V)$.
For the holomorphic degree-one cochain $a=0,b=1$, it is
$\det^{-1/2}$ in their oscillator convention. The reciprocal
half-determinant in a covariant Taylor-value convention is a dual
representation convention, not a change that may be silently made
on one side of an identity. On $\SU(V)$ the determinant twists
are trivial. As in 032, Lemma 4.4, pp.20--21, the circle and double
lifts agree there: their ratio is a continuous character of a
connected semisimple real group and hence is trivial.

At the compact places the appropriately split Gaussian is the
trivial representation; the exact possibilities are Liu,
Lemma D.2(1), pp.127--128. At the active place the holomorphic
component is $\omega_{m-1,1}^{-1,-,0}$ by Lemma D.2(2).
Fix a source splitting $\nu_m$ with
$\nu_m|_{\A_F^\times}=\epsilon_{E/F}^m$. Changing $\nu_m$
twists the source $\U(1)$ character, denoted by $\pi$; it does not
change the indicated $\mu$-component after that twist. In the
van Hoften--Sempliner convention the lift integrates
$\theta(g,t)\pi(t)$. Writing it instead as
$\theta(g,t)\eta(t)^{-1}$ requires $\eta=\pi^{-1}$.

At fixed finite level and prescribed infinity type the source
torus quotient has only finitely many character components.
Thus the raw first theta gradient is a finite sum of the corresponding
integrated lifts, all of whose nonzero components satisfy
\eqref{A6-extraction}. The degree-one argument of
van Hoften--Sempliner, Proposition 4.1.8, applies in every rank:
BMM, Section 9.1 and Theorem 9.1, pp.72--75, require
$3(a+b)+|a-b|<2m$, which here is $4<2m$. This theorem is being
used in degree one, not for the subsequent top-degree wedge.

For a finite list of forms, construct the full $\U(V)$ lifts with
these splittings and restrict to a connected component. A sufficiently
deep principal $\SU(V)$ level is contained in their common finite
stabilizer. Equivalently, at a lattice-preserving unitary level the
determinants of rational arithmetic elements are norm-one algebraic
units, hence roots of unity; a deep principal level kills them.
This gives the required component cover and its rational Hodge
pullbacks. No extension of an arbitrary unrealized special-unitary
form is being assumed.

There is also a split-local check on the rank parameter. Liu,
Lemma D.1, proof p.126, gives, for
$\mu=\nu\boxtimes\nu^{-1}$,
\[
 \omega(\mu,\varepsilon,\chi)
 =\operatorname{Ind}_{\mathrm{GL}_{m-1}\times\mathrm{GL}_1}
                       ^{\mathrm{GL}_m}
       ((\nu\circ\det)\boxtimes\chi\nu^{1-m}),
\]
with normalized induction. Its projective unramified Satake tuple is
\[
 \{q^{(m-2)/2},q^{(m-4)/2},\ldots,q^{-(m-2)/2},
                   \chi(\varpi)\nu(\varpi)^{-m}\}.
\]
Thus the exponent is $-m$, agreeing with $-3$ in 032. An inverse
geometric Hecke convention inverts the tuple. This is a consistency
check; Theorems 4.15 and 4.18 supply the actual CM source.

\paragraph{Degree $n$ alternation on the $n$-ball.}
Now set $m=n+1$. Choose multiplicity lines $P_1,\ldots,P_n$ with
the first $n$ prescribed signs, and put $P=\perp_{i=1}^nP_i$.
At the active place all line signs agree, and
\[
 \mathcal F_P=(L\otimes P)\oplus(Q^*\otimes P^*),
 \qquad \operatorname{rank}Q=n.
\]
The determinant computation in A4 gives
\[
 \det\mathcal F_P=L^{2n}(\det V_\mu)^{-n}(\det P)^{2-m}
\]
times compact constant factors. Its square root is $L^n$ times a
constant line. Define
\[
 \operatorname{Alt}_n:\operatorname{Sym}^n(Q^*\otimes P^*)
       \longrightarrow\bigwedge^nQ^*\otimes\bigwedge^nP^*
\]
by the product of the two exterior wedges. Swapping two factors
changes both wedge signs, so it factors through the symmetric power.
After the fixed constant-line trivializations, its target with the
half-determinant is
$L^n\otimes\bigwedge^nQ^*=\Omega_D^n$.

For a decomposable finite input the completed theta series is the
product of its $n$ line theta series. Its degree-$n$ coefficient,
in the divided-power convention, is the sum indexed by
$k_1+\cdots+k_n=n$. A summand with some $k_i\geq2$ contains
two factors in the same multiplicity line and is killed by
$\bigwedge^nP^*$. The sole surviving index is $(1,\ldots,1)$.
Consequently the projected degree-$n$ tensor is exactly the wedge
of the $n$ first tensors, with coefficient one.

For an arbitrary finite input the same assertion holds as a finite
sum of such wedges. Indeed a compactly supported locally constant
function on a finite-adelic direct sum is a finite sum of rectangle
indicators after choosing a sufficiently fine open lattice, so its
Schwartz space is the algebraic tensor product of the factor spaces.
A change to a rational adapted polarization uses the rational
metaplectic covariance of 032, Lemma 4.2 and Remark 4.3, in the
corresponding dimension. Its nonzero lift scalar is retained in the
constant-line normalization; it is not discarded outside $\SU(V)$.

Here is the spanning argument in the new dimension. Fix a point
$x_*\in D$ and, for one prescribed multiplicity line, choose rational
Schr\"odinger positions $F^m$. Their Minkowski image is a real vector
space of dimension $D_0=m[F:\Q]$, with full lattice
$\Lambda=\OO_F^m$. In the graph frame of its period point
$Z=(Z_\tau)_\tau$, put
\[
 P_i:\C^{D_0}\longrightarrow
       (L\otimes Q^*)_{x_*}.
\]
This is the complex-linear projection of the first Taylor fibre
onto its $c\mu$ summand, followed by the fixed nonzero
half-determinant and multiplicity-line trivializations. It is
surjective. The complex span of $\Lambda$ is $\C^{D_0}$, so choose
$x_{i,1},\ldots,x_{i,n}\in\Lambda$ such that
$P_i x_{i,1},\ldots,P_i x_{i,n}$ is a basis of the cotangent fibre.
For each $x_0=x_{i,j}$ use the actual finite input
$\phi_{x_0,k}=1_{x_0+k\widehat{\OO}_F^{\,m}}$. Its rational
support is exactly $x_0+k\OO_F^m$, and its projected first tensor is
\[
 2\pi i\sum_{x\in x_0+k\Lambda}
       P_i x\,
       \exp\left(\pi i\sum_\tau x_\tau^tZ_\tau x_\tau\right).
\]
Choose $\lambda_0>0$ below the least eigenvalue of every
$\operatorname{Im}Z_\tau$, and let $\rho>0$ be a lower bound for
the norm of a nonzero lattice vector. For
$k\geq2\|x_0\|/\rho$ and $0\ne y\in\Lambda$,
$\|x_0+ky\|\geq k\|y\|/2$. The tail is therefore bounded by
\[
 2\|P_i\|\sqrt{\frac{\pi}{e\lambda_0}}
   \sum_{r\geq1}(2r+3)^{D_0}
      \exp\left(-\frac{\pi\lambda_0 k^2\rho^2 r^2}{8}\right).
\]
Indeed $u e^{-\pi\lambda_0u^2/2}\leq
(\pi e\lambda_0)^{-1/2}$ absorbs the linear factor, and disjoint
balls of radius $\rho/2$ give at most $(2r+3)^{D_0}$ lattice
points in the shell $r\rho\leq\|y\|<(r+1)\rho$.
The displayed series tends to zero by dominated convergence;
THETA3-Q2 also gives a finite certified tail bound. Thus the
projected tensors tend to
\[
 2\pi i\,P_i x_{i,j}\,
 \exp\left(\pi i\sum_\tau
              x_{i,j,\tau}^tZ_\tau x_{i,j,\tau}\right).
\]
Each exponential is nonzero. The determinant of these $n$ limit
columns is nonzero, so continuity of the determinant gives $n$
actual finite inputs whose values form a basis. This construction
works separately for every prescribed type, with its own positive
Gaussian and projection $P_i$.

Now choose one input from each of the first $n$ prescribed families
in order. The already chosen values span a proper subspace until
the last step, while the next family's selected finite basis spans
the full fibre. One of its values lies outside that subspace.
The resulting $n$ values are independent. Their first theta forms
are holomorphic because the first completed tensor is the ordinary
holomorphic gradient and the active embedding projection is
holomorphic. Their wedge $\xi$ is nonzero. Only after this finite
selection do we choose a common principal level, using smoothness
of the finite Weil action. On the resulting smooth projective
compact ball quotient $X$,
\begin{equation}\label{A6-positive}
 i^{n^2}\int_X\xi\wedge\bar\xi>0.
\end{equation}
The phase is necessary when $n$ is odd: in complex coordinates,
$i^{n^2}dz_1\wedge\cdots\wedge dz_n\wedge
d\bar z_1\wedge\cdots\wedge d\bar z_n=2^n\,d\operatorname{vol}$.

\paragraph{A nearby rational orthogonal decomposition.}
Keep $V$, $P$, its rational symplectic realization, the finite total
input, and the total special-unitary level fixed. At a compact real
place, \eqref{A6-balance} says that the first and last lists contain
the same number of each sign. Relabel an initial real orthogonal
basis to have the desired last-list labels. This leaves its aggregate
positive and negative multiplicity subspaces unchanged. At the
active place every line has the same sign, so the total active
positive Lagrangian is unchanged as well.

Approximate these ordered orthogonal bases at all real places by
vectors in $P(E)$, using weak approximation. Apply Hermitian
Gram--Schmidt. The denominators are nonzero norms; being nonzero
with the required signs is an open condition. The resulting
rational orthogonal basis can therefore be chosen arbitrarily close
and with all the prescribed signs. It gives
$P=P_{n+1}\perp\cdots\perp P_{2n}$ over $E$.

Positive Gaussian majorants, with polynomial factors for derivatives,
give $C^\infty$ continuity of the projected total tensor on a compact
parameter neighborhood, exactly as in 032, Proposition 4.6(iii),
pp.22--24. Choose its constant-line trivializations continuously,
agreeing with the original ones at the reference point. The mixed
pairing with $\xi$ remains nonzero near that point by
\eqref{A6-positive}. Intermediate parameters need only give smooth
forms. At the chosen rational endpoint, finite tensor decomposition
gives a finite sum of wedges of holomorphic first tensors of the last
$n$ types. Select a summand whose mixed period is nonzero. After
a common principal cover for the selected inputs, one has
\begin{equation}\label{A6-period}
 \int_X\theta_1\wedge\cdots\wedge\theta_n
       \wedge\bar\theta_{n+1}\wedge\cdots\wedge\bar\theta_{2n}
        \ne0.
\end{equation}
The cover multiplies the integral by its positive degree.

The full unitary splitting characters may also be matched globally.
Choose $\mu_1,\ldots,\mu_{2n-1}$ with their prescribed types and set
\[
 \mu_{2n}=\frac{\prod_{i=1}^n\mu_i}
                    {\prod_{i=n+1}^{2n-1}\mu_i}.
\]
Its restriction to $\A_F^\times$ is
$\epsilon_{E/F}^{n-(n-1)}=\epsilon_{E/F}$, and
\eqref{A6-balance} gives its required weight-one infinity type.
Thus $\prod_{i\leq n}\mu_i=\prod_{i>n}\mu_i$ globally, with the
same $\nu_m$ retained on the source side. The special-unitary
argument has already removed these determinant twists, but they
are available for a compatible full-unitary tensor rule.

\paragraph{The direct cycle extraction, including its sign.}
By \eqref{A6-extraction}, expand the finitely many classes in
\eqref{A6-period} in complex spans of rational Hodge images.
Conjugate the last $n$ expansions. Nonvanishing selects rational
Hodge maps
$h_i:U(v_i)\to H^1(X,\Q)$ and nonzero
$\alpha_i\in U(v_i)_{\lambda_i}$ with algebraic Betti
coefficients, whose integrated product is nonzero. The arbitrary
analytic scalars are coefficients of the expansion; each selected
eigenvector itself is chosen over $\overline\Q$.

By the Albanese and the rational Riemann theorem, realize suitable
positive integer multiples of the $h_i$ by morphisms
$f_i:X\to B(v_i)$, as in 032, Lemma 2.2, p.5. Put
$A=\prod_iB(v_i)$, $g=\dim A$, $f=(f_i)$, and
$Z=\operatorname{cl}(f_*[X])\in H^{2g-2n}(A,\Q)$.
Writing $t_\alpha=\bigotimes_i\alpha_i$ in the distinct
$H^1$ slots, the projection formula gives
\[
 \int_A Z\smile t_\alpha=k\in\overline\Q^\times.
\]
For each $i$ take a copied factor $B(v_i)'$ and a nonzero
algebraic-Betti vector $\beta_i$ in its line labelled by
$c\lambda_i$. The pure class
$D_i=\alpha_i\otimes\beta_i\in H^2(B_i\times B_i',\overline\Q)$
is algebraic. Indeed every scalar conjugate puts its two factors
in opposite Hodge types, by centrality of $c$; scalar descent and
the Lefschetz $(1,1)$ theorem give its exact class as an algebraic
linear combination of divisor classes. No projection of an
unspecified divisor is involved.

Let $A'=\prod_i B_i'$, let $p,p'$ be the two projections from
$A\times A'$, and let $q_i$ select $B_i\times B_i'$. With the old
factors ordered before the new ones,
\[
 K:=\prod_iq_i^*D_i=(-1)^n t_\alpha\otimes t_\beta,
 \qquad t_\beta=\bigotimes_i\beta_i.
\]
The number of odd-degree transpositions is
$(2n)(2n-1)/2=n(2n-1)$, of parity $n$. Hence
\begin{equation}\label{A6-cycle}
 p'_*(p^*Z\smile K)=(-1)^n k\,t_\beta.
\end{equation}
Its codimension is $(g-n)+2n-g=n$. All operations on the left
preserve algebraicity, and $k$ is nonzero algebraic. Divide by $k$
and apply scalar complex conjugation, which preserves the
$\overline\Q$ span of rational cycle classes and changes the labels
back from $(c^n,1^n)$ to $(1^n,c^n)$. Identifying the copied factors
proves algebraicity of \eqref{A6-line}. This is a cycle on the line,
with the codimension and coefficient given by \eqref{A6-cycle}.
\hfill$\square$

\paragraph{The Grothendieck framework.}
This is an explicit composition of divisor correspondences,
pullback and proper pushforward in homological motives. The
relevant framework locators in the supplied Grothendieck scans
are \emph{cote 44}, scans 7--8, for intermediate Jacobians and
polarizations through correspondences, and \emph{cote 16},
scans 43--46, for the dimensions of Lefschetz operators and
algebraic inverse correspondences. Here the actual correspondence
$K$ is displayed and its output is computed in
\eqref{A6-cycle}. No general algebraic inverse Lefschetz
operator or Chow projector is assumed. Its action on the
specified tensor line follows directly from the divisor classes
and the projection formula.

\paragraph{An explicit family not generated by divisor products.}
For every $n\geq2$, choose a Galois CM field with at least $n+1$
real pairs of embeddings; a sufficiently large cyclotomic field
does so. Choose compact-place representatives
$\tau_1,\ldots,\tau_n$. Let $v_0$ be $+1$ on a chosen set of
representatives, including the distinguished one, and extend it
oddly. For $1\leq i\leq n$, let $v_i$ differ from $v_0$ only at
$\tau_i$ and its conjugate. Let $v_{n+1}$ differ at all those
$n$ pairs and put $v_{n+2}=\cdots=v_{2n}=v_0$. Then
\[
 \sum_{i=1}^n v_i=v_{n+1}+(n-1)v_0.
\]
No first-list function equals a second-list function.

The selected line \eqref{A6-line} is not in the subalgebra generated
by divisor classes. To prove this using the full Galois criterion,
project a product of $n$ divisor classes to that pure
one-$H^1$-per-factor line. A potentially nonzero contribution must
pair the $2n$ factors. At the distinguished embedding, each pair
must join a first-list factor to a last-list factor. For its selected
two eigenlines to have type $(1,1)$ under every scalar conjugation,
the two complete type functions must be equal. There are no such
pairs. Every divisor product therefore has zero projection to the
line. This argument uses the fixed eigenline labels, not a claim
that the underlying abelian varieties are pairwise nonisogenous.
Galois translates or isogenies involving other labels do not change
the conclusion. Formula \eqref{A6-cycle} nevertheless constructs
the line by a codimension-$n$ algebraic cycle, for all $n$, including
$n\geq4$.

\paragraph{Transfer to actual Weil pieces.}
Let a CM abelian variety $A_0$ of dimension $2n$ have multiplication
by an imaginary quadratic field $K$ with signature $(n,n)$.
Each of its two Weil lines over $\overline\Q$ is a wedge of
$2n$ CM embedding lines and has exactly $n$ of each Hodge type
under every scalar conjugation. Choose a Galois CM field containing
the CM fields involved, and let $s_i(g)$ be the complete type
functions of those lines. Reorder by their type at the distinguished
embedding and put $v_i=s_i$ for the first $n$, $v_i=-s_i$ for the
last $n$. Then \eqref{A6-balance} holds, and the selected original
line has the same complete type as $U(v_i)_{\lambda_i}$.

Extend a map between these matching lines by zero on the other
source embedding lines. Every scalar conjugate is a Hodge map,
so scalar descent expresses it as a $\overline\Q$ linear
combination of rational Hodge maps; this is 032, Lemma 2.4, p.8.
Lemma 2.2 realizes those maps by pullbacks of abelian morphisms.
Apply their tensor product to \eqref{A6-line} and take cup product.
After expansion this is a linear combination of pullbacks along
product morphisms from $A_0$ to the auxiliary product, and hence
preserves algebraicity. It proves the prescribed Weil line
algebraic, without an assumed algebraic projector.

For an explicit family in every $n$, choose a totally real Galois
field $F$ of degree $2n$, put $E=F(i)$, and select a CM type
containing $n$ embeddings above each embedding of $\Q(i)$.
The CM abelian variety has dimension $2n$ and the required Weil
signature. One field construction takes a prime
$\ell\equiv1\pmod{4n}$, the unique degree-$2n$ subfield of
$\Q(\zeta_\ell)$, which is real, and then adjoins $i$.
This does not assert simplicity or that every such Weil class is
outside the divisor algebra.

The same argument applies to the embedding-line monomials of any
rational CM Hodge class: rationality makes every occurring monomial
balanced under all scalar conjugations. Thus the direct
$2n$-factor construction can replace the four-factor combinatorial
reduction, with Liu's all-rank theorem retained as a deep cited input.

The remaining requirements for the proposed iterated Bessel bridge
are different. They include an actual tuple satisfying the recovered
\texttt{LawsX} and link conditions, identification of these
cohomological vectors and their normalizers, and the relevant
integral theta frames. Increasing $m$ does not satisfy the
every-place signature hypotheses of a different tower theorem,
and an independent $n$-fold product is not thereby an iteration
of theta operators.
